A lo largo de este texto se han presentado conceptos básicos sobre criptografía, los cuales han servido para contextualizar el algoritmo de AES. Además, se ha ido construyendo la base matemática necesaria para comprender qué son los cuerpos finitos y cómo operar con sus elementos. De esta forma, se puede entender cómo funciona cada paso de AES. 

Específicamente, la función \textit{SubBytes} y la función \textit{MixColumns}, donde se hace uso de estructuras algebraicas como el cuerpo finito \begin{Revision}$\Z_2[x]/(x^8 + x^4 + x^3 + x + 1)$\end{Revision} o el anillo $\F_{2^8}[x]/(\texttt{01}x^4+\texttt{01})$. El principal uso de \textit{SubBytes} es generar una función invertible no lineal que genere confusión al encriptar. En \textit{MixColumns} se toma una función invertible lineal para generar difusión, multiplicando las columnas del \textit{estado} por un elemento de dicho anillo que es coprimo con $\texttt{01}x^4+\texttt{01}$, y usando la identidad de Bézout se concluye que dicha función tiene inversa.

Debido a las limitaciones de este trabajo, no se han estudiado en profundidad otros conceptos relacionados con AES, como pueden ser los modos de operación, su proceso de diseño o el estudio de su criptoanálisis. Sin embargo, este trabajo proporciona una base sólida para el posterior estudio de estos temas más avanzados.